% Part: first-order-logic
% Chapter: proof-systems
% Section: axiomatic-deduction

\documentclass[../../../include/open-logic-section]{subfiles}

\begin{document}

\iftag{FOL}
      {\olfileid[hi]{fol}{prf}{axd}}
      {\olfileid[hi]{pl}{prf}{axd}}

\olsection{अभिगृहीतीय \usetoken{p}{derivation}}

सबसे पुराने और सरल तार्किक !!{derivation}-तंत्र अभिगृहीतीय
!!{derivation}-तंत्र हैं; उनकी !!{derivation}s केवल !!{sentence}s के क्रम हैं।
!!{sentence}s का क्रम सही !!{derivation} तभी है जब हर !!{sentence}~$!A$:
\begin{enumerate}
\item $!A$ अभिगृहीत हो, अथवा
\item $!A$ दिए गए समुच्चय~$\Gamma$ का !!{element} हो, जहाँ वह !!{sentence}s
  का समुच्चय है, अथवा
\item $!A$ किसी अनुमान-नियम से औचित्ययुक्त हो।
\end{enumerate}
अभिगृहीत $!A$ किसी नियत !!{sentence}-प्रतिरूप के रूप का हो। ऐसे प्रतिरूपों के
कई समुच्चय प्रथम-कोटि तर्कशास्त्र के लिए संतोषजनक (सुदृढ़ और पूर्ण)
!!{derivation}-तंत्र देते हैं। कुछ संयोजकों के अनुसार व्यवस्थित हैं; जैसे:
\[
!A \lif (!B \lif !A) \qquad !B \lif (!B \lor !C) \qquad (!B \land !C) \lif !B
\]
$\lif$, $\lor$ और~$\land$ के सामान्य अभिगृहीत हैं। कुछ तंत्र उनकी संख्या
घटाते हैं; आदिम संयोजकों के चुनाव पर एक अभिगृहीत भी पर्याप्त हो सकता है।

अनुमान-नियम किसी !!{sentence} को !!{derivation} में औचित्य देने की पर्याप्त
शर्त वाला आपादन है। मोडस पोनेन्स में $!A$ और $!A \lif
!B$ औचित्ययुक्त हों तो $!B$ भी। अतः !!{derivation} की !!{sentence}~$!B$
वाली पंक्ति तभी औचित्ययुक्त है जब $!A$ और $!A \lif !B$ दोनों (किसी
!!{sentence}~$!A$ के लिए) उस !!{derivation} में~$!B$ से पहले हों।

अभिगृहीतीय !!{derivation}-आधारित $\Proves$ संबंध में $\Gamma \Proves !A$ तभी
और केवल तभी जब ऐसी !!{derivation} हो जिसका अंतिम सूत्र !!{sentence}~$!A$ हो;
$\Gamma$ उन !!{sentence}s का समुच्चय है जिन्हें उस !!{derivation} में (2) से
औचित्य मिला। $!A$ तब प्रमेय है जब~$!A$ की !!{derivation} में~$\Gamma$ रिक्त हो,
यानी हर !!{sentence} उस !!{derivation} में (1) अथवा~(3) से औचित्ययुक्त हो।
मसलन, यह !!{derivation} दिखाती है कि $\Proves
!A  \lif (!B \lif (!B \lor !A))$:
\begin{derivation}
  1. & $!B \lif (!B \lor !A)$ \\
  2. & $(!B \lif (!B \lor !A)) \lif (!A  \lif (!B \lif (!B \lor !A)))$\\
  3. & $!A  \lif (!B \lif (!B \lor !A))$
\end{derivation}
पंक्ति~1 का !!{sentence}, अभिगृहीत $!A \lif (!A
\lor !B)$ का रूप है (यहाँ $!A$ और $!B$ की भूमिकाएँ उलटी हैं); पंक्ति~2 का
वाक्य अभिगृहीत $!A \lif (!B \lif !A)$ का रूप है। अतः दोनों औचित्ययुक्त हैं।
पंक्ति~3 को मोडस पोनेन्स से औचित्य मिलता है: इसे $!D$ लिखें तो पंक्ति~2 का रूप
$!C \lif !D$ है, जहाँ $!C$ का मान $!B \lif (!B \lor !A)$ अर्थात् पंक्ति~1 है।

समुच्चय $\Gamma$ असंगत है यदि $\Gamma \Proves \lfalse$। पूर्ण अभिगृहीत-तंत्र
$\lfalse \lif !A$ को हर~$!A$ के लिए सिद्ध करता है; अतः असंगत $\Gamma$ से
$\Gamma \Proves !A$ हर~$!A$ के लिए।

तर्कशास्त्र की अभिगृहीतीय !!{derivation}s पहले गॉटलोब फ़्रेगे की 1879 की
\emph{बेग्रिफ़्सश्रिफ़्ट} में मिलीं; इसीलिए वह आधुनिक तर्कशास्त्र की पहली कृति
मानी जाती है। इन्हें व्हाइटहेड–रसेल की \emph{प्रिंसिपिया मैथेमैटिका} और 1920
के दशक में हिल्बर्ट–शिष्यों ने परिष्कृत किया; इसलिए नाम ``फ़्रेगे-तंत्र'' या
``हिल्बर्ट-तंत्र'' भी हैं। ऐसा तंत्र प्रायः आसानी से मिलता है और
!!{derivation} की सरल संरचना तथा एक-दो अनुमान-नियम उसके \emph{विषय में}
परिणाम सिद्ध करना आसान बनाते हैं; पर व्यवहार में प्रमाण ढूँढ़ना और लिखना कठिन
है।

\end{document}
